% GENERATED FILE. Edit canonical lessons, then run npm run generate:books.
% canonical-source-hash: fnv1a64:4508d297a3cf3181
% canonical-lessons: LA-C552-domicilium, LA-C552-conclave, LA-C552-supellex, LA-C552-verro, LA-C552-scopae, LA-R552-review-545-552

\chapter{Rooms and furniture}
\label{ch:la-rooms-and-furniture}
\hlchaptermodality{552}

\begin{chapteropening}
\textbf{By the end of this chapter:} \emph{I can talk about rooms, furniture and sweeping.}

Complete the last lesson of chapter 552: I can talk about rooms, furniture and sweeping.
\end{chapteropening}

\section[domicilium, domiciliī]{domicilium, domiciliī — a home, dwelling}

\label{lesson:LA-C552-domicilium}



\begin{quote}
Before the new one: say the Latin for writing tablets, then the Latin for to write back.
\end{quote}

\subsection*{You'll want to know: domicilium, domiciliī}
\textbf{domicilium, domiciliī} — \textquotedblleft{}a home, dwelling\textquotedblright{}.

\begin{quote}
Haec vīlla domicilium familiae nostrae est.
\end{quote}

\begin{quote}
\emph{This villa is our family's home.}
\end{quote}

Neuter, second declension, from \emph{domus}. English \emph{domicile}.

\begin{cousinweb}
\textbf{domicilium} is built on \emph{domus} \textquotedblleft{}house\textquotedblright{}; Lewis \& Short sees in its second part the root of \emph{cēlō} \textquotedblleft{}to hide\textquotedblright{}. English \emph{domicile} comes from it.
\end{cousinweb}

\begin{grammarlens}[title={What you've built}]
One more word for telling, explaining and describing at B1.
\end{grammarlens}

\subsection*{Your turn}
\begin{itemize}
\raggedright
  \item \emph{Say it:} \emph{domicilium, domiciliī}
  \item \emph{Say it:} \emph{domicilium, domiciliī}, once more
  \item \emph{Say it:} a sentence of your own with \emph{domicilium, domiciliī}
  \item \emph{Recall:} say the Latin for writing tablets, then the Latin for to write back, then say \emph{domicilium, domiciliī} again
\end{itemize}

\begin{tcolorbox}[breakable,colback=teal!4,colframe=teal!35!black,title={Before you move on}]
What does domicilium, domiciliī mean? (\textquotedblleft{}A home, dwelling\textquotedblright{}.) Say it once more.
\end{tcolorbox}
\section[conclāve, conclāvis]{conclāve, conclāvis — a room}

\label{lesson:LA-C552-conclave}



\begin{quote}
Before the new one: say the Latin for to write back, then the Latin for a home.
\end{quote}

\subsection*{You'll want to know: conclāve, conclāvis}
\textbf{conclāve, conclāvis} — \textquotedblleft{}a room\textquotedblright{}.

\begin{quote}
Hoc conclāve parvum sed lūcidum est.
\end{quote}

\begin{quote}
\emph{This room is small but bright.}
\end{quote}

Neuter, third declension: a room that can be locked with one key, \emph{clāvis}. English \emph{conclave}.

\begin{cousinweb}
\textbf{conclāve} is built on \emph{con-} and \emph{clāvis} \textquotedblleft{}key\textquotedblright{} (Lewis \& Short): a room shut with one key. English \emph{conclave} comes from it.
\end{cousinweb}

\begin{grammarlens}[title={What you've built}]
One more word for telling, explaining and describing at B1.
\end{grammarlens}

\subsection*{Your turn}
\begin{itemize}
\raggedright
  \item \emph{Say it:} \emph{conclāve, conclāvis}
  \item \emph{Say it:} \emph{conclāve, conclāvis}, once more
  \item \emph{Say it:} a sentence of your own with \emph{conclāve, conclāvis}
  \item \emph{Recall:} say the Latin for to write back, then the Latin for a home, then say \emph{conclāve, conclāvis} again
\end{itemize}

\begin{tcolorbox}[breakable,colback=teal!4,colframe=teal!35!black,title={Before you move on}]
What does conclāve, conclāvis mean? (\textquotedblleft{}A room\textquotedblright{}.) Say it once more.
\end{tcolorbox}
\section[supellex, supellectilis]{supellex, supellectilis — household furniture, goods}

\label{lesson:LA-C552-supellex}



\begin{quote}
Before the new one: say the Latin for a home, then the Latin for a room.
\end{quote}

\subsection*{You'll want to know: supellex, supellectilis}
\textbf{supellex, supellectilis} — \textquotedblleft{}household furniture, goods\textquotedblright{}.

\begin{quote}
Supellex nostra simplex est: lectī, mēnsa, sellae.
\end{quote}

\begin{quote}
\emph{Our furniture is simple: beds, a table, chairs.}
\end{quote}

Feminine, third declension, singular only: all the movable things of a house. Ablative \emph{supellectile}.

\begin{cousinweb}
\textbf{supellex} is usually traced, as in Lewis \& Short, to \emph{super} and \emph{legō} \textquotedblleft{}to gather\textquotedblright{}, but the origin is not settled.
\end{cousinweb}

\begin{grammarlens}[title={What you've built}]
One more word for telling, explaining and describing at B1.
\end{grammarlens}

\subsection*{Your turn}
\begin{itemize}
\raggedright
  \item \emph{Say it:} \emph{supellex, supellectilis}
  \item \emph{Say it:} \emph{supellex, supellectilis}, once more
  \item \emph{Say it:} a sentence of your own with \emph{supellex, supellectilis}
  \item \emph{Recall:} say the Latin for a home, then the Latin for a room, then say \emph{supellex, supellectilis} again
\end{itemize}

\begin{tcolorbox}[breakable,colback=teal!4,colframe=teal!35!black,title={Before you move on}]
What does supellex, supellectilis mean? (\textquotedblleft{}Household furniture, goods\textquotedblright{}.) Say it once more.
\end{tcolorbox}
\section[verrō, verrere]{verrō, verrere — to sweep}

\label{lesson:LA-C552-verro}



\begin{quote}
Before the new one: say the Latin for a room, then the Latin for household furniture.
\end{quote}

\subsection*{You'll want to know: verrō, verrere}
\textbf{verrō, verrere} — \textquotedblleft{}to sweep\textquotedblright{}.

\begin{quote}
Ancilla pavīmentum ātriī verrit.
\end{quote}

\begin{quote}
\emph{The maid sweeps the floor of the hall.}
\end{quote}

\emph{verrō, verrere, verrī, versum}.

\begin{cousinweb}
\textbf{verrō} is inherited from a Proto-Indo-European root meaning \textquotedblleft{}to sweep, drag along the ground\textquotedblright{} (EtymDB).
\end{cousinweb}

\begin{grammarlens}[title={What you've built}]
One more word for telling, explaining and describing at B1.
\end{grammarlens}

\subsection*{Your turn}
\begin{itemize}
\raggedright
  \item \emph{Say it:} \emph{verrō, verrere}
  \item \emph{Say it:} \emph{verrō, verrere}, once more
  \item \emph{Say it:} a sentence of your own with \emph{verrō, verrere}
  \item \emph{Recall:} say the Latin for a room, then the Latin for household furniture, then say \emph{verrō, verrere} again
\end{itemize}

\begin{tcolorbox}[breakable,colback=teal!4,colframe=teal!35!black,title={Before you move on}]
What does verrō, verrere mean? (\textquotedblleft{}To sweep\textquotedblright{}.) Say it once more.
\end{tcolorbox}
\section[scōpae, scōpārum]{scōpae, scōpārum — a broom}

\label{lesson:LA-C552-scopae}



\begin{quote}
Before the new one: say the Latin for household furniture, then the Latin for to sweep.
\end{quote}

\subsection*{You'll want to know: scōpae, scōpārum}
\textbf{scōpae, scōpārum} — \textquotedblleft{}a broom\textquotedblright{}.

\begin{quote}
Servus scōpīs culīnam pūrgat.
\end{quote}

\begin{quote}
\emph{The slave cleans the kitchen with a broom.}
\end{quote}

Feminine, first declension, plural: a bundle of twigs. Italian \emph{scopa} — broom — comes from it.

\begin{cousinweb}
\textbf{scōpae} first meant \textquotedblleft{}twigs\textquotedblright{} (Lewis \& Short), and a bundle of twigs made a broom. Lewis \& Short connects it with a root meaning \textquotedblleft{}to prop\textquotedblright{}, but the origin is not settled.
\end{cousinweb}

\begin{grammarlens}[title={What you've built}]
One more word for telling, explaining and describing at B1.
\end{grammarlens}

\subsection*{Your turn}
\begin{itemize}
\raggedright
  \item \emph{Say it:} \emph{scōpae, scōpārum}
  \item \emph{Say it:} \emph{scōpae, scōpārum}, once more
  \item \emph{Say it:} a sentence of your own with \emph{scōpae, scōpārum}
  \item \emph{Recall:} say the Latin for household furniture, then the Latin for to sweep, then say \emph{scōpae, scōpārum} again
\end{itemize}

\begin{tcolorbox}[breakable,colback=teal!4,colframe=teal!35!black,title={Before you move on}]
What does scōpae, scōpārum mean? (\textquotedblleft{}A broom\textquotedblright{}.) Say it once more.
\end{tcolorbox}
\section[(dialogue)]{Words from chapters 545–552 — a review}

\label{lesson:LA-R552-review-545-552}



\begin{quote}
Say the Latin for a sound and to sound.
\end{quote}

\subsection*{Your turn}
Say these aloud:

\begin{itemize}
\raggedright
  \item a sound, to sound, to resound, a band of musicians, harmony
  \item the seats of a theatre, a stage, a theatre curtain, a prologue, a mime dancer
  \item an image, a likeness, a shape, a small cube, to engrave
  \item beauty, good-looking, elegant, an ornament, ugly
  \item a poem, an epigram, a little verse, a public reading, a listener
  \item learned, a bookseller, a reader, to read through, to unroll
  \item a book box, a writing case, parchment, writing tablets, to write back
  \item a home, a room, household furniture, to sweep, a broom
\end{itemize}

\begin{tcolorbox}[breakable,colback=teal!4,colframe=teal!35!black,title={Before you move on}]
A sound? (\textbf{sonus, sonī}.) To sound? (\textbf{sonō, sonāre}.) To resound? (\textbf{resonō, resonāre}.) A band of musicians? (\textbf{symphōnia, symphōniae}.) Harmony? (\textbf{concentus, concentūs}.) The seats of a theatre? (\textbf{cavea, caveae}.) A stage? (\textbf{pulpitum, pulpitī}.) A theatre curtain? (\textbf{aulaeum, aulaeī}.) A prologue? (\textbf{prologus, prologī}.) A mime dancer? (\textbf{pantomīmus, pantomīmī}.) An image? (\textbf{imāgō, imāginis}.) A likeness? (\textbf{effigiēs, effigiēī}.) A shape? (\textbf{figūra, figūrae}.) A small cube? (\textbf{tessella, tessellae}.) To engrave? (\textbf{caelō, caelāre}.) Beauty? (\textbf{pulchritūdō, pulchritūdinis}.) Good-looking? (\textbf{speciōsus, speciōsa, speciōsum}.) Elegant? (\textbf{ēlegāns, ēlegantis}.) An ornament? (\textbf{ōrnāmentum, ōrnāmentī}.) Ugly? (\textbf{dēfōrmis, dēfōrme}.) A poem? (\textbf{poēma, poēmatis}.) An epigram? (\textbf{epigramma, epigrammatis}.) A little verse? (\textbf{versiculus, versiculī}.) A public reading? (\textbf{recitātiō, recitātiōnis}.) A listener? (\textbf{audītor, audītōris}.) Learned? (\textbf{litterātus, litterāta, litterātum}.) A bookseller? (\textbf{bibliopōla, bibliopōlae}.) A reader? (\textbf{lēctor, lēctōris}.) To read through? (\textbf{perlegō, perlegere}.) To unroll? (\textbf{ēvolvō, ēvolvere}.) A book box? (\textbf{capsa, capsae}.) A writing case? (\textbf{scrīnium, scrīniī}.) Parchment? (\textbf{membrāna, membrānae}.) Writing tablets? (\textbf{pugillārēs, pugillārium}.) To write back? (\textbf{rescrībō, rescrībere}.) A home? (\textbf{domicilium, domiciliī}.) A room? (\textbf{conclāve, conclāvis}.) Household furniture? (\textbf{supellex, supellectilis}.) To sweep? (\textbf{verrō, verrere}.) A broom? (\textbf{scōpae, scōpārum}.)
\end{tcolorbox}
